\subsection{Relative Topology}

\begin{definition} \label{def:relative-open-subset} \hlt{Relative Openness}\\
Suppose $E\subseteq Y\subseteq X$, where $X$ is a metric space. The set $E$ is \hlt{open relative to $Y$} if for each $p\in E$ there is an $r>0$ such that $q\in E$ whenever $q\in Y$ and $d(p,q)<r$.
\end{definition}

\begin{theorem} \label{thm:relative-open-intersection-with-open-set} \hlt{Relative Opens Are Intersections}\\
Suppose $Y\subseteq X$. A subset $E$ of $Y$ is open relative to $Y$ if and only if $E=Y\cap G$ for some open subset $G$ of $X$.
\begin{proof}
\begin{enumerate}[label=\roman*.]
\setlength{\itemsep}{0pt}
\item[$(\Rightarrow)$] For each $p\in E$, choose $r_p>0$ with $B_{r_p}(p)\cap Y\subseteq E$, and set $G\defeq\bigcup\limits_{p\in E}B_{r_p}(p)$. Then $G$ is open in $X$, and $E\subseteq G\cap Y$. If $q\in G\cap Y$, then $q\in B_{r_p}(p)\cap Y$ for some $p\in E$, so $q\in E$. Thus $E=G\cap Y$.
\item[$(\Leftarrow)$] If $E=G\cap Y$ with $G$ open in $X$, then for each $p\in E$ choose $r>0$ with $B_r(p)\subseteq G$.\\
Hence $B_r(p)\cap Y\subseteq G\cap Y=E$, so $E$ is open relative to $Y$.
\end{enumerate}
\hlt{Structural Commentary}: Relative openness is ordinary openness viewed through the ambient subset $Y$.\\
\hlt{Key Insight}: Intersect open sets in $X$ with $Y$ to get exactly the open sets of $Y$.
\end{proof}
\end{theorem}

\begin{example} \label{ex:relative-topology-ambient-examples} \hlt{Relative Topology Depends on the Ambient Space}\\
Let $X=\{(x,0):x\in\R\}\subseteq\R^2$ with the metric inherited from $\R^2$. The set $E=\{(x,0):-1<x<1\}$ is not open as a subset of $\R^2$, but it is open relative to $X$, since $E=X\cap\bigl((-1,1)\times(-1,1)\bigr)$.\\
Likewise, $[0,1)$ is not closed in $\R$, but it is closed relative to $(-1,1)$ because $[0,1)=(-1,1)\cap[0,\infty)$.
\end{example}
